\chapter{Lumpectomy}
\section*{Overview}
Lumpectomy, or breast-conserving surgery, removes a breast tumor with a surrounding margin while preserving most of the breast.
\section*{Common Instruments and Equipment}
The exact setup varies by technique and setting. Common items may include:
\begin{itemize}[nosep,leftmargin=*]
  \item breast surgery instrument set
  \item localization wire or other lesion-localization device when needed
  \item specimen orientation markers and imaging equipment
\end{itemize}

\section*{Why It Is Done}
It is used for selected localized breast cancers and some high-risk or noncancerous lesions when breast conservation is appropriate.
\section*{How It Is Performed}
The lesion is localized when necessary, removed with a margin, and sent for pathology. Sentinel lymph-node biopsy may be performed during the same operation.
\section*{Contraindications}
No single contraindication applies to every patient; the treating team assesses whether the procedure is suitable for the individual and the reason for treatment.

\section*{Precautions}
The care team reviews relevant health conditions, medicines, allergies, and procedure-specific risks, and explains any changes needed for the individual.

\section*{Risks and Recovery}
Bleeding, infection, seroma, altered breast shape, positive margins requiring re-excision, and recurrence are possible. Radiation or systemic treatment may follow according to pathology.
\section*{References}
\begin{enumerate}\item National Cancer Institute. Breast-Conserving Surgery. 2025.\item Current Medical Diagnosis and Treatment. McGraw-Hill; 2025.\end{enumerate}
\clearpage
